% Part: computability
% Chapter: computability-theory
% Section: normal-form

\documentclass[../../../include/open-logic-section]{subfiles}

\begin{document}

\olfileid{cmp}{thy}{nfm}
\olsection{نارمل روپ دا قضیہ}

فرض کرو اسی قابلِ حساب فنکشناں دیاں تعریفاں بیان کر سکدے
آں، تے پرکھ سکدے آں کہ کسے انپٹ اُتے اوس فنکشن دا حساب کرن دے
پورے ریکارڈ دا دِتا ہویا ممکنہ بیان درست اے یا نہیں۔ فیر
ایہ گل معقول لگدی اے کہ حساب دے نمونے توں آزاد رہ کے، اسی
کسے وی قابلِ حساب فنکشن~$f$ دی کسے وی انپٹ~$x$ اُتے قدر
ایویں لبھ سکدے آں:
\begin{enumerate}
  \item حسابی ریکارڈاں دے سارے ممکنہ بیان تلاش کرو۔
  \item پرکھو کہ دِتا ریکارڈ $f(x)$ دے حساب دا ریکارڈ اے یا نہیں۔
  \item جے درست ریکارڈ لبھ جائے، تاں اوہدے وچوں $f(x)$ دی قدر کڈھو۔
\end{enumerate}
ایہ گل واقعی درست اے: کلینی دے نارمل روپ دے قضیے دا مضمون
ایہو اے۔

\begin{thm}[کلینی دے نارمل روپ دا قضیہ]
\ollabel{thm:normal-form}
اک بنیادی بازگشتی تعلق~$T(e, x, s)$ تے اک بنیادی بازگشتی
فنکشن~$U(s)$ موجود نیں، جیہناں دی ایہ خاصیت اے: جے $f$
کوئی وی جزوی قابلِ حساب فنکشن اے، تاں کسے~$e$ لئی
\[
f(x) \simeq U(\umin{s}{T(e, x, s)})
\]
ہر~$x$ اُتے۔
\end{thm}

\begin{proof}[ثبوت دا خاکہ]
حساب دے کسے وی نمونے وچ قابلِ حساب فنکشن~$f$ دا بیان پوری
ریاضیاتی سختی نال تعریف کیتا جا سکدا اے، تے ایس بیان نوں
اک قدرتی عدد~$e$ دا کوڈ دِتا جا سکدا اے۔ ``حسابی سلسلہ''
دا تصور وی پوری ریاضیاتی سختی نال تعریف کیتا جا سکدا اے،
جیہڑا اشاریہ~$e$ والے فنکشن دے انپٹ~$x$ اُتے حساب دا
عمل درج کردا اے۔ ایہوجیہے حسابی سلسلے نوں وی اک عدد~$s$
دا کوڈ دِتا جا سکدا اے۔ ایہ کم ایداں ہو سکدا اے کہ
\begin{enumerate}
  \item تعلق $T(e, x, s)$، جیہڑا صرف تے صرف اُدوں قائم ہوندا اے جدوں
  عدد~$s$، اشاریہ~$e$ والے فنکشن دے انپٹ~$x$ اُتے حسابی
  سلسلے دا کوڈ ہووے، تے
  \item فنکشن $U(s)$، جیہڑا کوڈ~$s$ والے حسابی سلسلے نوں
  اوس حساب دے آخری نتیجے ول نقشہ بند کردا اے،
\end{enumerate}
دونوں قابلِ حساب ہون۔ بلکہ تعلق~$T$ تے فنکشن~$U$ بنیادی
بازگشتی نیں۔
\end{proof}

\begin{explain}
نارمل روپ دے قضیے دا پوری ریاضیاتی سختی نال ثبوت دین لئی
سانوں حساب دا اک نمونہ مقرر کرنا، قابلِ حساب فنکشناں دے
بیاناں تے حسابی سلسلیاں دی کوڈنگ پوری تفصیل نال کرنا، تے
ایہ تصدیق کرنا پوے گا کہ تعلق~$T$ تے فنکشن~$U$ بنیادی
بازگشتی نیں۔ زیادہ تر اطلاقاں لئی ایہ کافی اے کہ $T$
تے~$U$ قابلِ حساب ہون تے $U$~کُل ہووے۔

نارمل روپ دے قضیے دا ثبوت شاید اک مختصر یاد رکھن والی گل
دے روپ وچ چیتے رکھنا سب توں چنگا اے: $\umin{s}{T(e, x, s)}$
اشاریہ~$e$ والے فنکشن دے انپٹ~$x$ اُتے حسابی سلسلے نوں
لبھدا اے، تے جے کوئی ایہوجیہا سلسلہ لبھ جائے تاں $U$
اوہدا نتیجہ واپس دیندا اے۔
\end{explain}

جے حساب دا نمونہ جزوی بازگشتی فنکشناں والا ہووے، جیہڑا
کلینی نے شروع وچ ورتیا سی، تاں ایہ نتیجہ دسّدا اے کہ
کسے وی فنکشن دی تعریف لئی غیر محدود تلاش صرف اک واری
ورتی جان دی لوڑ اے، یعنی اکّو $\umin{y}{f(\vec x, y)}$
آپریٹر کافی اے۔ ایس معنی وچ ایہ وکھاندا اے کہ ہر جزوی
بازگشتی فنکشن دا اک نارمل روپ اے۔

$T$ تے $U$ نال فہرست $\cfind{0}$، $\cfind{1}$، $\cfind{2}$،
\dots دی تعریف کیتی جا سکدی اے۔ ہن توں اسی فرض کراں گے
کہ $T$ تے~$U$ دی اک مناسب چُون مقرر کر لئی اے، تے مساوات
\[
\cfind{e}(x) \simeq U(\umin{s}{T(e,x,s)})
\]
نوں $\cfind{e}$ دی \emph{تعریف} منّاں گے۔

ایہ اک ہور کم دی حقیقت اے:

\begin{thm}
ہر جزوی قابلِ حساب فنکشن دے لامتناہی اشاریے نیں۔
\end{thm}

ایہ گل وی سوجھ نال صاف سمجھ آؤندی اے۔ کسے قابلِ حساب
فنکشن دا بیان ملے، تاں اک ہور بیان بنایا جا سکدا اے
جیہڑا اوہو فنکشن، یعنی اوہو انپٹ--آؤٹ پٹ جوڑا، حسابے،
پر تھوڑے وکھرے طریقے نال۔ مثلاً اوہ پہلے کوئی ایہوجیہا
کم کرے جیہدا حساب اُتے کوئی اثر نہ ہووے: پرکھے کہ
$0 = 0$ اے، یا $5$ تک گنے، وغیرہ۔ بدلے ہوئے بیان دا
اشاریہ اصل اشاریے توں ہمیشہ وکھرا ہووے گا۔ دونوں اوہو
فنکشن دے اشاریے نیں؛ بس اوہدا حساب تھوڑے وکھرے طریقے
نال کیتا گیا اے۔

\end{document}
